\chapter{ECSA Exit Level Outcomes}

This section addresses the following Engineering Council of South Africa (ECSA) exit level outcomes and indicates how each outcome was demonstrated in the project.

\section*{GA 1 - Problem-Solving}

The project addresses an engineering problem which constitutes the lack of a joint design for Portable Coordinate Measurement Arms (PCMAs) that provides a repeatabile, reliable position-holding capability. This was solved by designing, modelling, and validating an electromagnetic brake joint capable of producing a sufficient holding torque.

\section*{GA 2 - Application Scientific and Engineering Knowledge}

To satisfy this outcome, fundamental principles across mathematics, science, and computing must be applied to create effective solutions for complex engineering problems. This was demonstrated by using mathematical formulae to determine the required magnetic flux density, force generation, and torque produced to ensure the design achieved the sufficient torque. Principles of mechanics and physical science were used to reason and conceptualise the friction produced, the load-bearing capabilities, and the locking mechanics necessary to achieve reliable position-holding.

\section*{GA 3 - Engineering Design}

This outcome requires the design of creative solutions for complex engineering problems and the design of systems, components, or processes to meet an identified need. This was achieved by executing a satisfactory design process that included a requirements analysis, concept generation and evaluation between pnematic and electromagnetic locking, as well as mechanical friction methods. The design of the chosen electromagnetic brake concept was refined and modelled in detail, and complemented by a dedicated switching circuit. The design was finalised following a techno-analysis covering workshop fabrication costs and component availability.

\section*{GA 5 - Use of Engineering Tools}
 
This outcome requires the demonstration of competence in creating, selecting, and applying appropriate techniques, resources, and modern engineering or IT tools. Autodesk Inventor was used for 3D CAD configuration, component packaging, and modelling the assembly of the PCMA joint. LTSpice was used to simulate and optimise the switching ciruit driving the electromagnet. Experimental test apparatus and measurements instruments were used to validate physical holding torque, while accounting for simulation idealisations and sensor limits.

\section*{GA 6 - Professional and Technical Communication}

This outcome requires demonstrating the competence to communicate effectively on complex engineering activities, both orally and in writing. Written communication was demonstrated through a structured report adhering to departmental standards, supported by 3D CAD mechanical models, electronic circuit schematics, and empirical performance graphs. Oral and visual communication was demonstrated through the progress presentation, final oral presentation, and poster presentation designed for a general technical audience.

\section*{GA 8 - Individual and Collaborative Teamwork}

This outcome requires demonstrating competence to function effectively as an indiviual. Although execution encountered significant schedule delays, individual competence was demonstarted by actively managing these setbacks, re-evaluating critical path dependencies, and strategically prioritising core project objectives. Work was systematically refocused onto primary deliverables, specifically the design, driving circuit, and experimental validation of the electromagnetic joint. This was done while maintaining communication with the supervisor and coordinating with departmental workshop staff and laboratory technicians to bring the project to completion.

\section*{GA 9 - Independent Learning Ability}

This outcome requires engaging in independent learning through self-directed research skills. Independent literature reviews were conducted electromagnetic actuator design, high-current switching circuits, and joint stiffness optimisation. These were topics that extended beyond standard coursework. This self-acquired knowledge was synthesised and directly applied to design, build, and evaluate the final joint prototype.
